\documentclass[10pt,journal,compsoc]{IEEEtran}

\usepackage{amsmath,amssymb,amsfonts}
\usepackage{mathrsfs}
\usepackage{booktabs}
\usepackage{url}
\newcommand{\botrule}{\bottomrule}

\begin{document}

\title{Supplementary Material for\\ SENTRY: Anytime-Valid E-Process Monitoring
of LLM Agent Action Streams}

\author{Quang-Vinh~Dang%
\thanks{Q.-V. Dang is with British University Vietnam, Hung Yen, Vietnam.
E-mail: vinh.dq4@buv.edu.vn.}}

\markboth{IEEE Transactions on Dependable and Secure Computing --- Supplementary Material}%
{Dang: SENTRY (Supplementary Material)}

\maketitle

\noindent This supplement collects the reproducibility settings, the extended
signal definitions, and the Shiryaev--Roberts versus CUSUM discussion referred
to in the main paper. Section, equation and reference numbers below are
self-contained; equations reused from the main paper are restated here for
convenience. All code, the collected trajectory corpus, and the scripts that
regenerate every reported number and figure are released with the project
repository.

\section{Reproducibility and Hyperparameters}\label{app:repro}

Every number and figure in the paper is regenerated by scripts released with
the project. The evaluation is fully deterministic: feature hashing uses a
seed-stable CRC rather than Python's per-process-salted \texttt{hash}, so
back-to-back runs of the evaluation produce byte-identical output. Randomness
enters only through the ten train/calibration/threshold/test splits, which are
seeded \(0,\dots,9\).

\begin{table}[h]
\centering
\caption{Hyperparameters. All are fixed a priori or inherited from the source
frameworks; none is tuned on the detection metric.}\label{tab:hyper}
\begin{tabular}{@{}lll@{}}
\toprule
Symbol & Meaning & Value \\
\midrule
\(\alpha\) & target false-alarm rate & \(0.2\) \\
\(\delta\) & PAC confidence parameter & \(0.2\) \\
\(\lambda_{\mathrm{nov}}\) & novelty weight & \(4\) \\
\(\lambda_{\mathrm{IL}}\) & instruction weight & \(4\) \\
\((w_1,\dots,w_4)\) & instruction-likeness weights & \((3,2,2,1)\) \\
\(\kappa\) & power-calibrator exponent & \(0.5\) \\
\(K\) & tilt-grid size (exp.\ mixture) & \(6\) \\
splits & random evaluation splits & \(10\) (seeds \(0\)--\(9\)) \\
min actions & trajectory length filter & \(2\) \\
\botrule
\end{tabular}
\end{table}

The reference model is fit only on the training split of \emph{nominal}
trajectories; no attack trajectory is seen during fitting or calibration. Data
collection used DeepSeek-V4-Flash (served through OpenRouter) as the agent
(and, for \(\tau\)-bench, the simulated-user) backend; the AgentDojo pipeline
used the \texttt{important\_instructions} attack, and the change point of each
attack was located by matching the injected string against observation content
after alphanumeric normalisation.

\section{Extended Signal Definitions}\label{app:signals}

\paragraph*{Observation instruction-likeness.} Writing \(o_{t-1}\) for the tool
observation preceding action \(a_t\), the raw instruction-likeness
\(\mathrm{IL}(o_{t-1})\) is a per-word weighted density of generic
instruction cues, and the shipped signal is the excess over a self-calibrating
nominal ceiling \(\beta\),
\begin{equation}
s^{\mathrm{IL}}_t = \max\!\big(0,\; \mathrm{IL}(o_{t-1}) - \beta\big),
\qquad
\beta = \max_{o \in \mathcal{D}_{\mathrm{train}}} \mathrm{IL}(o).
\label{eq:il}
\end{equation}
The lexicon is fixed and generic: imperative verbs (\emph{please, send, ignore,
forward, do, must, before, after, execute, run, \dots}), second-person
pronouns (\emph{you, your, yours}), pseudo-markup blocks matched by
\verb|<[A-Za-z/][^>]{0,40}>|, and exclamation marks; densities are per word.
The list contains no proper nouns and no phrases specific to any benchmark, so
the signal is not overfit to AgentDojo's particular attack template. Because
\(\beta\) is the maximum instruction-likeness observed on the training split,
the excess in~\eqref{eq:il} is zero on any observation no more directive than
nominal tool output and positive only on genuinely anomalous,
instruction-bearing observations.

\paragraph*{Argument novelty and tool transition.} The argument-novelty
vocabulary is maintained per tool, with a global fallback, over tokens matched
by \verb|[A-Za-z0-9_./@-]{2,}|. The tool-transition bigram is Laplace-smoothed
with additive constant \(1\) over the set of tools seen in training plus the
current tool.

\section{On the Choice of Shiryaev--Roberts vs.\ CUSUM}\label{app:srcusum}

SENTRY ships both e-detector recursions. Writing \(L_n\) for the per-action
baseline increment and \(M_n\) for the detector statistic, the
Shiryaev--Roberts (SR) and CUSUM recursions are
\begin{align}
\text{SR:}\quad & M_n = L_n\,(M_{n-1} + 1), & M_0 &= 0, \label{eq:sr}\\
\text{CUSUM:}\quad & M_n = L_n\,\max(M_{n-1},\, 1), & M_0 &= 1. \label{eq:cusum}
\end{align}
We report SR throughout because, in the mixture-of-increments setting, its
additive accumulation (\(M_{n-1}+1\)) responds slightly faster to the short
attacks in our corpus than CUSUM's \(\max(M_{n-1},1)\); the two are known to
have matching first-order asymptotic delay under the e-detector framework
[Shin, Ramdas and Rinaldo, \emph{New England Journal of Statistics in Data
Science}, 2023], and on our data the difference in detection rate is within
the cross-split standard deviation.

\end{document}
